分级风险、切面准确率和弱目标重叠分别评价严重程度预测、采集切面识别及与局部标注的一致程度。完整模型取得最低分级风险，而删除变体具有更高的空间任务指标，说明评价时应同时考察评分及各项伴随输出。模块删除比较评价的是整体设计，未分离监督、模型容量和信息使用方式的影响，任务取舍背后的机制仍待明确。

局部矩形框为空间学习提供位置线索。将区域响应图与原始图像及医生标注框对照，可以直接检查响应与这些线索的对应关系，并发现可供进一步查看的区域。弱目标重叠衡量位置一致性，完整病灶覆盖则是另一项评价目标。框外响应可能对应相关组织或错误激活，更完整的空间参考与医生复核有助于区分这些情况。

对于临床复核，联合输出提供三个检查对象：预测切面、区域响应和最终评分。所示病例中评分接近但切面预测错误的情况，展示了这些输出如何呈现同一预测的不同方面。将它们与原始图像并列呈现，可以分别检查解剖背景和局部响应。其实际获益仍有待通过阅片者研究评价，具体包括错误识别、判断修订的恰当性和复核时间。

部署研究应在本文报告的模型开销之外，于目标硬件上联合测量预处理、推理和结果展示。独立队列评价应检验不同患者人群与采集条件下的分级及空间输出表现。
